\subsection{A four-element opposite corner on four states}

\begin{lemma}[Proved: a corner retract retaining a nonzero restriction]
\label{CP2:A-corner-retract-nonzero-restriction}
Let $C$ be a two-object collage with corners $M,N$ and
nonempty connecting biset, and suppose $BN$ is rationally
acyclic. If $BC\simeq\CP^2$, and an actual categorical
retraction $C\to C'$ has an inclusion section which retains
the entire corner $M$, then $BC'\simeq\CP^2$.
\end{lemma}
\begin{proof}
The rational generator $u\in H^2(BC;\mathbb Q)$ restricts
nontrivially to $BM$. Otherwise it vanishes on both corners
and lifts from their actual suspension quotient in the
two-block model, so $u^2=0$. Its restriction to $BC'$ is
therefore nonzero. The retraction makes the integral
homology of $BC'$ a direct summand of that of $BC$, and
its fundamental group a retract of the trivial group.
Consequently $H_2(BC';\mathbb Z)=\mathbb Z$ and its split
degree-two map is an isomorphism. Naturality of the primitive
square forces the split degree-four map to be an isomorphism
as well. All other positive homology groups vanish. The
inclusion is thus an integral homology equivalence between
simply connected CW spaces, and hence a homotopy equivalence.
\end{proof}

\begin{lemma}[Proved: the nilpotent four-corner extension through order twelve]
\label{CP2:A-nilpotent-four-four-states-twelve}
Let $N=\{1,e_0,e_1,a\}$ satisfy
\[
 e_i e_j=e_j,\qquad ae_i=e_i,\qquad e_i a=e_0,
 \qquad a^2=e_0.
\]
For a nonempty left-$N$, right-$M$ biset $S$ with
$|S|\le4$ and $|M|\le12$, its collage cannot have the
rational homology and cohomology ring of $\CP^2$.
\end{lemma}
\begin{proof}
We spell out where the source bound enters the complete
comparison of Proposition~\ref{CP2:A-nilpotent-four-corner-four-states}.
Write $P_0,P_1,F$ for the left actors and $U$ for the common
image of $P_0,P_1$. These retractions fix $U$ pointwise,
$F$ fixes $U$, and every right actor preserves $U$. If
$|U|=1$, this joint fixed state supplies a section of the
bridge projection to $BM$; all second classes are relative
to the corners and have zero products. Suppose $|U|\ge2$.

Put $R=\mathbb QM$, $L_U=\ker(\mathbb QU\to\mathbb Q)$,
$D=(1-P_0)\mathbb QS$ and $H=F-P_0$. The full exact
right-projective resolution of the quoted proposition gives
the natural reduced bridge-projection complex
\[
 C_0=L_U,\qquad C_j=D\ (j\ge1),\qquad
 d_1=P_1-P_0,\quad d_j=H\ (j\ge2).
\]
Here $H^2=d_1H=0$, and the actual collage, after collapsing
its rationally acyclic corner $BN$, has reduced chains
$(C\otimes_R^{\mathbf L}\mathbb Q)[1]$. The inclusion
$L_U\to C$ is the comparison induced by the actual
sub-biset $U$, not just an abstract map with the same ranks.

Since $\dim D\le2$, either $H=0$, or $\dim D=2$ and
$\ker H=\operatorname{im}H$ is one-dimensional. In the
first case the exact direct sum
\[
 C=[D\xrightarrow{d_1}L_U]\ \oplus\
       \bigoplus_{j\ge2}D[j]
\]
shows that bounded rational collage homology forces
$D\otimes_R^{\mathbf L}\mathbb Q$ to be acyclic.
Thus the actual subcollage on $U$ includes by a rational
homology equivalence. On it the $N$ action is trivial;
collapsing $BN$ gives an actual rational homology and ring
equivalence to the ordinary $M$-act cone on $U$. Boundedness
passes through these specified equivalences. Its products
vanish by Theorem~\ref{CP2:A-bounded-twelve-acts}.

In the rank-one case the natural quotient of $C$ is the
two-term complex
\[
 [D/\operatorname{im}H\xrightarrow{\overline d_1}L_U]
\]
with both terms one-dimensional. If its differential is
nonzero, it is an isomorphism of right $R$-modules, so the
collage is rationally acyclic. Otherwise $P_0=P_1$ as
functions and $P_0:S\to U$ is an actual proper biset
retraction. As proved in the quoted proposition, this
retraction preserves a rational projective-plane homology
ring: the relative quotient is a suspension, so its image
must have nonzero $H^2$, and the split maps then give the
same rational homology and nonzero square. The image has
trivial left action and is rationally equivalent to the
ordinary $M$-act cone on $U$. It is bounded, so the same
bounded-twelve theorem excludes it. These are all cases;
the source order was used only at this ordinary-act step.
\end{proof}

\begin{proposition}[Proved: the complete nineteen-arrow budget $(11,4,4)$]
\label{CP2:A-oriented-nineteen-eleven-four-four}
A two-object collage with corner orders $11,4$ and four
connecting arrows cannot have type $\CP^2$.
\end{proposition}
\begin{proof}
Orient the connecting arrows from the order-eleven monoid
$M$ to the order-four monoid $N$, with left-$N$, right-$M$
biset $S$. Opposition covers the reversed orientation.
If the trivial right $\mathbb QN$ module is projective,
the actual orbit quotient is a rational homology and ring
equivalence to an ordinary $M$-act cone. Its homology is
bounded under the proposed hypothesis, so
Theorem~\ref{CP2:A-bounded-twelve-acts} excludes it.

The full minimal-ideal classification in
Theorem~\ref{CP2:A-six-versus-four-semisimple-quotient}
leaves exactly six nonprojective order-four corners. One
is the nilpotent monoid of the preceding lemma. The other
five are the identity-adjoined right-zero band
\[
 R_4=\{1,e_0,e_1,e_2\},\qquad e_i e_j=e_j,
\]
and the four monoids $\{1,e_0,e_1,a\}$ with
\[
 e_i e_j=e_j,\quad ae_i=e_i,\quad
 e_i a=e_{\tau(i)},\qquad
 (\tau,a^2)=(\mathrm{id},1),(\mathrm{id},a),
 ((01),1),(\mathrm{const}_0,a).
\]
Each has a right zero $e_i$, and hence contractible $BN$.
Write $P_i$ for its band actors, $A$ for the extra actor
when present, and $U$ for the common image of all $P_i$.
They fix $U$ pointwise; the extra actor also fixes $U$.
All right actors preserve $U$ and commute with these maps.

We may suppose $|U|\ge2$: a unique core state would be
jointly fixed and its bridge section would kill the source
restriction and hence every square. If $S=U$, the whole
left action is trivial. The actual extreme-corner retract
of Theorem~\ref{CP2:A-trivial-extreme-corner-retract} gives
a smaller projective-plane category, contrary to the
completed eighteen-arrow exclusion
\ref{CP2:A-global-lower-nineteen}. We will use that exclusion
for every subsequent proper retract. A proper equivariant
state retraction has projective-plane image by
Theorem~\ref{CP2:A-cut-retract}. A proper corner retraction
which preserves $M$ has the same consequence by
Lemma~\ref{CP2:A-corner-retract-nonzero-restriction}.

First take $N=R_4$. If at most two different maps occur
among $P_0,P_1,P_2$, choose one band element for each
different map. Their union with $1$ is a proper submonoid.
Sending every band element to its chosen representative
is a monoid retraction, is compatible with the literal
action on $S$, and fixes the source $M$. It therefore gives
the forbidden smaller collage. Suppose the three maps
are distinct. If $|U|=3$, their values at the one remaining
state are all three core states. The left act is the
regular left $R_4$-act, whose action category has an initial
object. Thus the bridge projection to $BM$ is a rational
homology equivalence and the collage is rationally acyclic.

Only $|U|=2$, with two noncore states, remains. Label the
core $u_0,u_1$. Each $P_i$ is a two-bit word on the two
noncore states, and three distinct such words occur.
The two column profiles in $\{0,1\}^3$ are therefore
nonconstant, different, and not complementary. A commuting
map restricts to a self-map of $U$. If it is the identity
there, the two distinct profiles force it to fix both
other states. If it is constant on $U$, the nonconstant
profiles force the same global constant map. If it swaps
$U$, the complementary profiles that it would require
do not occur, so this is impossible. The full commuting
transformation monoid is consequently exactly
\[
 \{\operatorname{id}_S,c_{u_0},c_{u_1}\}.
\]
The right $M$-image either fixes a core state, giving the
joint fixed-state obstruction, or is this entire monoid.
The latter is excluded by
Theorem~\ref{CP2:A-two-constant-image}.

For the remaining four corners, if $P_0=P_1=P$ then
$P$ commutes with $A$ and all right actors. Since $U\ne S$,
it is a proper biset retraction. Thus assume $P_0\ne P_1$.
For $(\mathrm{id},a)$ the map $A$ is a central idempotent
actor. Unless it is the identity it supplies a proper
biset retraction. If it is the identity, the corner
retraction $a\mapsto1$ onto $R_3=\{1,e_0,e_1\}$ is
compatible with $S$ and yields the forbidden smaller
collage.

For $(\mathrm{id},1)$, if $A$ is the identity the same
corner retraction works. Otherwise $A$ is an involution
fixing $U$ pointwise. Since $2\le|U|<|S|\le4$, it must
swap exactly two noncore states, with $|U|=2$. The maps
$P_i$ are constant on this orbit. Since they differ,
their two endpoint values are different. The natural
actual graph in Lemma~\ref{CP2:A-four-corner-actual-graph}
has one midpoint and its two legs join the two core
vertices: it is a tree.

For $((01),1)$ the equation $P_iA=P_{1-i}$ rules out
$A=1$ under $P_0\ne P_1$. Again $A$ swaps the two
noncore states and fixes the two core states. The same
graph has one midpoint with different endpoints, hence
is a tree.

Finally take $(\mathrm{const}_0,a)$. Here $A^2=A$ and
$P_iA=P_0$. If $A=P_0$, the corner retraction
$a\mapsto e_0$ onto $R_3$ respects the action. Otherwise
$AS$ strictly contains $U$: if $AS=U$ then
$A=P_0A=P_0$. It is also properly contained in $S$,
since an idempotent surjection is the identity, which
would force $P_0=P_1$. Therefore $|S|=4$, $|U|=2$ and
$|AS|=3$. On $AS$ the maps $P_0,P_1$ agree, so their
distinct values occur at the sole state outside $AS$.
The actual graph in the cited lemma again has exactly
one midpoint, with different endpoints, and is a tree.

In all three tree cases the graph comparison is a natural
actual rational homology equivalence. Its augmented
chains are acyclic, and the nonnegative right-$M$ bar
therefore makes the bridge projection to $BM$ a rational
homology equivalence. Since $BN$ is rationally acyclic,
the full collage is rationally acyclic. These cases cannot
have projective-plane type. Together with the nilpotent
lemma they exhaust all order-four corners.
\end{proof}

\begin{remark*}[Scope of the small-action check]
The companion literal check enumerates every left action
of the six displayed four-element monoids on one through
four labelled states, and verifies the asserted retract,
regular-act, centralizer or tree alternative. It is a
check of the stated finite case split. The proof uses no
catalog of eleven-element source monoids and no completion
claim about any larger action search.
\end{remark*}
